% !TeX root = ../navier-stokes-manuscript.tex
% ============================================================
% 3.3 The Critical Height at the Terminal Face
% ============================================================
\subsection{The Critical Height at the Terminal Face}\label{sub:TheCriticalHeightAtTheTerminalFace}

The terminal rectangles provide adjacent spatial rows.  Their first task is to
turn two time-integrated norms into one strong value at the terminal face.

\begin{lemma}[Adjacent Hilbert-scale trace]\label{lem:AdjacentHilbertScaleTrace}
Let \(s\in\R\) and suppose
\[
  v\in L^2(I_*;H^{s+1}),
  \qquad
  v_t\in L^2(I_*;H^{s-1}).
\]
Then \(v\) has one representative in \(C([0,T_*];H^s)\), and for
\(0\leq a\leq b\leq T_*\),
\begin{equation}\label{eq:AdjacentHilbertScaleEstimate}
  \left|
  \norm{v(b)}_{H^s}^{2}
  -\norm{v(a)}_{H^s}^{2}
  \right|
  \leq
  2\norm{v_t}_{L^2(a,b;H^{s-1})}
  \norm{v}_{L^2(a,b;H^{s+1})}.
\end{equation}
\end{lemma}

\begin{proof}
Let \(A=(1-\Delta)^{1/2}\).  Time mollification and frequency truncation make
the following identity legitimate for smooth approximants:
\begin{equation}\label{eq:AdjacentHilbertScaleIdentity}
  \norm{v(b)}_{H^s}^{2}
  -\norm{v(a)}_{H^s}^{2}
  =
  2\operatorname{Re}\int_a^b
  \pair{A^{s-1}v_t(t)}{A^{s+1}v(t)}_{L^2}
  \,dt.
\end{equation}
Cauchy--Schwarz gives \eqref{eq:AdjacentHilbertScaleEstimate}.  The two assumed
\(L^2\) rows pass the right-hand side to the limit.  The limit identity makes
the \(H^s\) norm continuous, while the approximation also gives weak
\(H^s\) continuity; together they imply strong continuity.  Any two continuous
representatives agree everywhere because they agree almost everywhere.
\end{proof}

The critical readout needs only \(\cR_{0,2}[u_0]\).  Its integer rows embed into
the fractional adjacent pair
\begin{equation}\label{eq:SectionThreeCriticalAdjacentPair}
  u\in L^2(I_*;H^{5/2}),
  \qquad
  u_t\in L^2(I_*;H^{1/2}).
\end{equation}
No fractional rectangle has been assumed; these are continuous embeddings of
the integer record.  The trace lemma at \(s=3/2\) now gives
\begin{equation}\label{eq:SectionThreeCriticalVelocityTrace}
  u\in C([0,T_*];H^{3/2}).
\end{equation}

Let \(H=E_{1/2}\).  To read the homogeneous critical energies from the
inhomogeneous trace, mollify in time and truncate frequencies to
\(\varepsilon\leq|\xi|\leq R\).  Differentiation and Plancherel then give
\[
  \frac{1}{2}\norm{\Lambda^{3/2}u(b)}_{L^2}^{2}
  -\frac{1}{2}\norm{\Lambda^{3/2}u(a)}_{L^2}^{2}
  =
  \operatorname{Re}\int_a^b
  \pair{\Lambda^{1/2}u_t(t)}{\Lambda^{5/2}u(t)}_{L^2}
  \,dt.
\]
Cauchy--Schwarz passes the integral to the limit through the two critical rows;
the continuous \(H^{3/2}\) trace passes the endpoint terms.  Removing all
cutoffs yields, for \(0\leq a\leq b\leq T_*\),
\[
  E_{3/2}(b)-E_{3/2}(a)
  =
  \operatorname{Re}\int_a^b
  \pair{\Lambda^{1/2}u_t(t)}{\Lambda^{5/2}u(t)}_{L^2}
  \,dt.
\]
The adjacent pair makes the derivative integrable, and the complete spatial row
makes \(E_{3/2}\) itself integrable.  Hence
\begin{equation}\label{eq:SectionThreeCriticalEnergyRegularity}
  E_{3/2}\in W^{1,1}(I_*)\cap L^1(I_*)
  \subset L^\infty(I_*).
\end{equation}

The integer rows also put the derivative of the critical height in \(L^1\):
\begin{equation}\label{eq:SectionThreeCriticalHeightDerivative}
  H'(t)
  =\pair{\Lambda^{1/2}u(t)}{\Lambda^{1/2}u_t(t)}_{L^2}
  \in L^1(I_*).
\end{equation}
Keep the pressure-explicit production current as
\[
  \mathcal P_{1/2}
  :=
  -\operatorname{Re}
  \pair{\Lambda^{1/2}u}
  {\Lambda^{1/2}\bigl((u\cdot\nabla)u+\nabla p\bigr)}_{L^2}.
\]
Pairing the full equation with \(\Lambda u\) gives
\[
  H'=-2\nu E_{3/2}+\mathcal P_{1/2}.
\]
The two integrable rows consequently imply
\[
  H\in W^{1,1}(I_*),
  \qquad
  E_{3/2}\in L^\infty(I_*)\cap L^1(I_*),
  \qquad
  \mathcal P_{1/2}\in L^1(I_*).
\]
The height and \(E_{3/2}\) therefore reach the terminal face continuously; the
production current is only integrable, and no endpoint value is claimed for it.
This entire critical trace comes from \(\cR_{0,2}[u_0]\).  Higher rectangles
will repeat the endpoint argument at every finite Sobolev height.
